\documentclass[11pt]{article}
\usepackage[a4paper,margin=26mm]{geometry}
\usepackage{amsmath,amssymb,amsthm,hyperref}
\newtheorem{theorem}{Theorem}
\newtheorem{lemma}{Lemma}
\newcommand{\PP}{\mathcal P}
\title{A conditional lifting theorem for several small dice}
\author{Complete proof; independently reviewed}
\date{30 September 2026}
\begin{document}\maketitle
\noindent\textbf{Status.} Full independent mathematical reviews by the root,
rational-design, and independent-direction agents passed. This is not formal
verification.

\noindent\textbf{Scope.} This is a conditional finiteization theorem.
Its hypotheses are verified by the rational smoothing construction in
\path{two_polynomial_exceptional_dice.tex}. For general $k$, the present
note is only a conditional criterion.

\begin{theorem}\label{main}
Let $m,k\ge1$. Suppose independent variables
$A_1,\ldots,A_k,U_1,\ldots,U_m$ are fully permutation-fair,
where each $U_j$ is uniform on $[0,1]$ and each $A_i$ is uniform on
$K_i$ distinct interior points. Suppose their atom locations admit a
partition into open intervals (clusters) such that:
\begin{enumerate}
\item all cluster closures are disjoint and contained in $(0,1)$;
\item each cluster contains atoms of exactly one $A_i$;
\item each cluster contains at least $m^2$ distinct atoms;
\item each cluster moment $K_i^{-1}\sum_{a\in C}a^r$ is rational
      for $0\le r\le m$.
\end{enumerate}
Then a completely finite permutation-fair family exists with exactly
$K_i$ faces on its $i$th specified die, $1\le i\le k$.
The other $m$ dice may be chosen to have one common finite face count.
No quantitative bound on that common count is asserted.
\end{theorem}

The point of the hypothesis is the number of distinct nodes in each
cluster. Replacing one atom by coalescing replicas does not justify it:
the Vandermonde inverse then deteriorates, and small geometric
perturbations need not yield small density corrections.

Fix a full order of all $m+k$ labels. Conditional on the $k$ specified
values $t_1,\ldots,t_k$, its probability for uniform old variables is
zero unless the specified labels have their prescribed relative order.
In the correct chamber it is
\begin{equation}\label{kernel}
 P(t)=\prod_{a=0}^k\frac{(t_{a+1}-t_a)^{r_a}}{r_a!},
 \qquad t_0=0,\quad t_{k+1}=1,\quad\sum_a r_a=m,
\end{equation}
where the specified variables have been listed in that order.
On a product of clusters its chamber is fixed, and $P$ is a polynomial
of total degree at most $m$ (or identically zero).

\begin{lemma}[Several boundary responses]\label{response}
Around each specified node $t_q$, choose a small interval, all such
intervals disjoint, and assign the node to at most one old label.
For a node assigned to label $j$, choose a signed density $h_q$ with
\[
 \int_{x<t_q}Q(x)h_q(x)\,dx=Q(t_q),\qquad
 \int_{x>t_q}Q(x)h_q(x)\,dx=-Q(t_q)
 \quad(Q\in\PP_{m-1}).
\]
Give old label $j$ density $1+\sum_{q\text{ assigned to }j}\delta_qh_q$.
For every specified variable $i$, cluster $C$, and old label $j$, suppose
\begin{equation}\label{common}
 \frac1{K_i}\sum_{q\in C:\ q\text{ assigned to }j}
                   \delta_q Q(t_q)=B_{i,C}(Q)
                 \qquad(Q\in\PP_{m-1}),
\end{equation}
with the same functional $B_{i,C}$ for every $j$.
Let $E_{i,C}=K_i^{-1}\sum_{q\in C}\operatorname{ev}_{t_q}$.
Then the probability of the fixed full order is
\begin{equation}\label{operator}
 \sum_{C_1,\ldots,C_k}
 \left[\prod_{i=1}^k(E_{i,C_i}+B_{i,C_i}\partial_i)\right]P.
\end{equation}
Here each operator acts only on its own coordinate, and the appropriate
polynomial branch of \eqref{kernel} is used on each cluster product.
\end{lemma}
\begin{proof}
Expand the product of the old densities in the order integral.
Each old variable contributes at most one bump. If the support of any
selected bump is not cut by one of the selected specified values, then,
after integrating the untouched uniform variables, its integrand is a
polynomial of degree at most $m-1$ throughout its whole support.
The support has fixed order relative to all other selected supports
and cut points. Its full moments vanish, so that term is zero.

Consequently a surviving term assigns its selected old variables
injectively to selected specified cut points; in particular there
are at most $k$ bumps. The truncated evaluation identities replace
each such old integration by evaluation at its cut, with a positive
sign on the lower side and a negative sign on the upper side.
They are precisely the boundary terms obtained by differentiating
\eqref{kernel} once in each of those cut coordinates.

For clarity about collisions, a gap containing $r$ prescribed old
labels contributes $l^r/r!$. A derivative at its lower endpoint
removes the first old label with a minus sign, and one at its upper
endpoint removes the last old label with a plus sign. If both act
and $r\ge2$, they remove distinct labels and give the corresponding
second derivative. If $r=1$, the second derivative is zero: there
is no term using the same old variable twice. If $r=0$, even the
first derivative is zero. Since each specified cut is differentiated
at most once, these are all possible collisions. Different gaps
factor independently. This proves the claimed mixed-derivative
interpretation with its exact coefficients.

Now sum over node choices for each fixed assignment of distinct
old labels to the cuts. On a cluster product all remaining expressions
are polynomials of degree at most $m-1$ in each differentiated
coordinate. Equation~\eqref{common} replaces every such sum by
$B_{i,C}$, independently of the old label that was assigned to it.
The undifferentiated coordinates give $E_{i,C}$. Summing over all
subsets of cut coordinates gives \eqref{operator}.
\end{proof}

\begin{proof}[Proof of Theorem~\ref{main}]
Choose rational cluster endpoints and disjoint rational support
intervals around all baseline nodes, with each support inside its
cluster. For every cluster choose $m$ disjoint groups of $m$ nodes,
one group for each old label. Unused nodes are unassigned.
Approximate all baseline nodes by distinct rational nodes $t_q$ inside
their fixed support intervals. Denote the original cluster functional
by $E^0_{i,C}$ and the new one by $E_{i,C}$.

Both have the same mass. Therefore
\begin{equation}\label{Bdef}
 B_{i,C}(Q)=E^0_{i,C}(JQ)-E_{i,C}(JQ),\qquad Q\in\PP_{m-1},
\end{equation}
is well defined, where $JQ$ is any polynomial antiderivative.
The rational cluster-moment hypothesis ensures that all $B_{i,C}$
are rational, although the baseline nodes themselves need not be.
They tend to zero as the rational approximations approach the
baseline nodes.

For each old label's group in each cluster, the evaluation map on
$\PP_{m-1}$ is an invertible Vandermonde matrix. Solve the rational
linear system \eqref{common} for the $\delta_q$. Because the original
nodes within every group are distinct, the inverse matrices stay
bounded in a sufficiently small neighborhood of the baseline.
Thus all $\delta_q$ tend to zero.

The required $h_q$ can be chosen rational and uniformly bounded as
$t_q$ varies near its baseline node. Choose $m$ disjoint rational
subintervals to the left of that neighborhood and $m$ to its right,
all inside the fixed support. The moment matrices of their indicator
functions are invertible: a polynomial of degree at most $m-1$
whose integral vanishes on each of $m$ disjoint intervals has a zero
in each interval, hence is zero. Solve separately for truncated
moments $Q\mapsto Q(t_q)$ and $Q\mapsto-Q(t_q)$. The fixed inverse
matrices and continuous right sides give the uniform bound.
Taking the rational approximation sufficiently close ensures that
all old densities stay positive, for instance between $1/2$ and $3/2$.

By construction
\[
 (E_{i,C}+B_{i,C}\partial_i)P=E^0_{i,C}P
 \quad(P\in\PP_m).
\]
The response lemma, successively in each coordinate on every cluster
product, reproduces exactly the original order probability. Hence the
new rational specified nodes and the positive rational step densities
of the old variables are fully fair.

Finally partition at every density breakpoint and every specified
node. All old interval masses are rational. Take any finite fair
$m$-die block with $F$ faces per die and a common denominator $D$
for those masses. Within each interval replace the continuous old
variables by this block, independently replicating each old letter
by its integer interval-mass numerator. Conditional on which old
variables fall in that interval, their order remains uniform.
Each old die has exactly $FD$ faces. Insert one specified face at
each of its rational nodes and assign globally distinct integer
labels. This preserves every complete order and finishes the proof
with the stated face counts.
\end{proof}

\noindent\textbf{Unresolved hypothesis.} For general $k$, this remains a conditional criterion. The companion
two-variable smoothing proof supplies the hypotheses when $k=2$.
Rationality of the cluster moments cannot simply be dropped when the
desired final model has equiprobable finite faces.
\end{document}
